\section{Motivation}
\label{sec:bg}

\paragraph{Blockchain databases.~\cite{ruan2021blockchains}.}
\XXX.
BlockchainDB~\cite{el2019blockchaindb}.
FalconDB~\cite{peng2020falcondb}.

Authenticated data structures~\cite{martel2004general, tamassia2003authenticated}.
\textsc{Integri}DB~\cite{zhang2015integridb}.

As the blockchain databases keep operating, the size of content they store may ever grow.
However, blockchains, as state machine replication protocols, by default require each node to store a full copy of the execution state in order to replicate the execution individually.
Eventually, the ever-growing database will overwhelm the storage capacity of the nodes. 

\paragraph{Sharded blockchains.}
To overcome this problem, the stored data can be divided into shards and each node stores only one or a subset of the shards.
However, without the whole execution state, the nodes can no longer execute every transaction.
One solution is to shard the historical ledger (\ie the blocks) and leave the execution state intact.
For example, BFT-Store~\cite{qi2020bft} performs erasure coding on the blocks to reduce the storage complexity per block to $O(1)$.
EIP-4844~\cite{eip4844} of Ethereum proposes to shard \emph{blob}, a special kind of block data only whose commitment is accessible by the transactions.
These works guarantee the availability of blocks and claim that one can always replay the execution state from an available chain.
However, online execution is still required by many blockchain systems to \emph{validate} the block.
Moreover, if the execution state size exceeds the capacity of a single node, even replaying becomes infeasible.

The other solution is to perform \emph{processing sharding} like regular databases -- each node processes a subset of the transactions that only make use of the state shard(s) it stores.
This is safe for regular databases, however, reducing the execution quorum immediately brings safety concerns to the blockchain databases, as the faulty nodes gain more power in a smaller execution quorum.
BlockchainDB~\cite{el2019blockchaindb}.
SharPer~\cite{amiri2021sharper}: ``the number of nodes, $N$, is assumed to be much larger than $f$''.
\xxx{Discuss randomize execution quorum approaches.}
Moreover, cross shard transactions bring additional coordination overhead to processing sharding approach, posing performance penalty and complicating the reasoning of safety.

% Efficient reliable broadcast with erasure code\cite{cachin2005asynchronous}, used in HoneyBadger\cite{miller2016honey}.
% Ceph~\cite{weil2006ceph}, HDFS~\cite{shvachko2010hadoop}, Azure~\cite{huang2012erasure}.
% Morph~\cite{kim2024morph}.

\paragraph{Decouple sharding from execution.}
To summarize, state sharding is desired for scaling out blockchains like databases, while global execution is desired for solid security.
However, this design space is rarely exploited.
Quote a survey of sharding replication systems from 2023~\cite{solat2023sharding}:
\emph{To the best of our knowledge, actually there is no protocol that uses storage sharding without processing sharding.}

Decoupling sharding from execution is challenging.
The storage layer must mask the missing shards prior to execution and synchronize them from storing nodes in time.
The synchronized shards must prove their \emph{integrity} \ie their contents are not tampered by the faulty storing nodes.
And most importantly, Byzantine faults must not cause any \emph{liveness} issues by withholding the storing shards.